% SOURCE_AUTHORITY: sga5_fr_workpass.tex, lines 13577--13612 (LF-normalized unit).
% SOURCE_SHA256: 791F4EFFC5E02832D5D77ED03518C8156D6F07E4C8238B03545DB93D883FBB28
Por último, $u(f'_R)\ne0$. En efecto, gracias a la condición b), existe un elemento $x_0\in R$ tal que $f'_R(x_0)-x_0\notin\fm^{r+1}$. Pero necesitamos un elemento de esta índole que pertenezca a $\fm$. Para obtenerlo, escribamos $x_0$ en la forma
\[
x_0=\frac{x}{x'},\qquad x,x'\in\fm,\quad x'\notin\fm^{r+1}.
\]
Al menos uno de los elementos $x,x'$ satisface la condición requerida.

Supongamos por reducción al absurdo que se tiene simultáneamente:
\[
f'_R(x) - x \in \fm^{r+1} \quad , \quad f'_R(x') - x' \in \fm^{r+1} \quad .
\]
Se tendrá entonces, en el cuerpo $K$ de fracciones de $R$:
\[
\frac{f'_R(x)}{x} \in 1 + \fm^r \quad , \quad \frac{f'_R(x')}{x'} \in 1 + \fm^r \quad .
\]
Se puede suponer que $f'_R(x') \neq 0$, pues, de lo contrario, $x'$ es el elemento buscado. Bajo esta hipótesis, teniendo en cuenta que $1 + \fm^r$ es un grupo multiplicativo, resulta que:
\[
\frac{x'}{f'_R(x')} \in 1 + \fm^r \quad ,
\]
y, por consiguiente:
\[
\frac{f'_R(x)}{x} \cdot \frac{x'}{f'_R(x')} \in 1 + \fm^r \quad .
\]
Pero esta última relación implica:
\[
\frac{f'_R(x_0)}{x_0} \in 1 + \fm^r \quad .
\]
En efecto:
\[
\begin{aligned}
\frac{f'_R(x_0)}{x_0} - \frac{f'_R(x)}{x} \cdot \frac{x'}{f'_R(x')}
&= \frac{f'_R(\frac{x}{x'})}{\frac{x}{x'}} - \frac{f'_R(x)}{x} \cdot \frac{x'}{f'_R(x')}
= \frac{x'\,f'_R(\frac{x}{x'})}{x} - \frac{f'_R(x)}{x} \cdot \frac{x'}{f'_R(x')} \\
&= \frac{1}{x}\left(\frac{x'\,f'_R(x')\,f'_R(\frac{x}{x'}) - x'\,f'_R(x)}{f'_R(x')}\right) = 0
\end{aligned}
\]
